\documentclass[10pt]{article}
\usepackage[T1]{fontenc}
\usepackage[utf8]{inputenc}
\usepackage[french]{babel}
\usepackage{amsmath,amssymb,mathrsfs,mathpazo}
\usepackage{geometry}
\geometry{a4paper,margin=2cm}
\usepackage{graphicx,array,tabularx,tabularray,colortbl,multirow,multicol,booktabs,enumitem,pifont,ulem,fancybox,diagbox,cancel}
\usepackage{pstricks,pstricks-add,pst-plot,pst-tree,pst-node,pst-text,pst-func,pst-math,pst-bspline,pst-3dplot,pst-eucl,pst-solides3d}
\let\pscalculate\undefined
\usepackage{pst-bezier}
\usepackage{tikz,pgfplots,tkz-tab}
\usetikzlibrary{arrows,arrows.meta,patterns,calc,positioning,trees,decorations.markings,shapes.geometric,shapes.misc,shapes.arrows,intersections,decorations.pathmorphing,decorations.pathreplacing,backgrounds,shadows,babel}
\pgfplotsset{compat=1.18}
\usepgfplotslibrary{fillbetween}
\usepackage[np]{numprint}
\usepackage{listings,makecell,pas-tableur,twoopt}
\let\barre\undefined
\usepackage[pstricks]{tablvar}
\definecolor{darkgreen}{RGB}{0,100,0}
\providecommand{\red}{\color{red}}
\providecommand{\blue}{\color{blue}}
\providecommand{\green}{\color{green}}
\providecommand{\black}{\color{black}}
\providecommand{\white}{\color{white}}
\providecommand{\gray}{\color{gray}}
\providecommand{\cyan}{\color{cyan}}
\providecommand{\magenta}{\color{magenta}}
\providecommand{\yellow}{\color{yellow}}
\providecommand{\violet}{\color{violet}}
\providecommand{\orange}{\color{orange}}
\providecommand{\euro}{\texteuro}
\providecommand{\EUR}{\texteuro}
\providecommand{\ieme}{\textsuperscript{e}}
\providecommand{\degres}{\ensuremath{^\circ}}
\providecommand{\degree}{\ensuremath{^\circ}}
\providecommand{\vv}[1]{\overrightarrow{#1}}
\providecommand{\Coord}[1]{#1}
\providecommand{\decosix}{}
\providecommand{\starredbullet}{\ensuremath{\bullet}}
\pagestyle{empty}
\setlength{\parindent}{0pt}
\usepackage{adjustbox}\DeclareUnicodeCharacter{2236}{:}\DeclareUnicodeCharacter{2190}{\ensuremath{\leftarrow}}
\providecommand{\R}{\mathbb{R}}
\providecommand{\N}{\mathbb{N}}
\providecommand{\Z}{\mathbb{Z}}
\providecommand{\Q}{\mathbb{Q}}
\providecommand{\D}{\mathbb{D}}
\providecommand{\C}{\mathbb{C}}
\providecommand{\e}{{\mathrm{e}}}
\providecommand{\vect}[1]{\overrightarrow{#1}}
\providecommand{\vectt}[1]{\overrightarrow{\mathrm{#1}}}
\providecommand{\vvt}[1]{\overrightarrow{\mathrm{#1}}}
\providecommand{\np}[2][]{{#2~#1}}
\providecommand{\vv}[1]{{\overrightarrow{#1}}}
\renewcommand{\d}{{\mathrm{d}}}
\providecommand{\degres}{{^\circ}}
\providecommand{\degre}{{^\circ}}
\providecommand{\blue}{\color{blue}}
\providecommand{\red}{\color{red}}
\providecommand{\enskip}{\quad}
\providecommand{\llbracket}{\lbrack\!\lbrack}
\providecommand{\rrbracket}{\rbrack\!\rbrack}
\providecommand{\euro}{{\text{€}}}
\providecommand{\no}{{\text{n°}}}
\providecommand{\interval}[2]{{[#1;#2]}}
\providecommand{\intervFO}[2]{{[#1;#2[}}
\providecommand{\dotfill}{\ldots}

\providecommand{\ds}{\displaystyle}
\providecommand{\textding}[1]{#1}
\providecommand{\Oij}{\(\mathrm{O};\vec{\imath},\vec{\jmath}\)}
\providecommand{\Oijk}{\(\mathrm{O};\vec{\imath},\vec{\jmath},\vec{k}\)}
\providecommand{\Ouv}{\(\mathrm{O};\vec{u},\vec{v}\)}
\providecommand{\intervFF}[2]{\ensuremath{[#1;#2]}}
\renewcommand{\euro}{\ifmmode\text{\texteuro}\else\texteuro\fi}
\DeclareUnicodeCharacter{00B0}{\ifmmode^{\circ}\else\textdegree\fi}
\DeclareUnicodeCharacter{20AC}{\ifmmode\text{\texteuro}\else\texteuro\fi}
% Énoncés : Éducation nationale. Transcription et corrigés : APMEP, Denis Vergès et contributeurs. Découpage et classement : CoopMaths, Sébastien Lozano, Éric Elter et contributeurs. Adaptation de la présentation : Maths Explorer.
% CC BY-SA : https://coopmaths.fr/www/about/
% https://coopmaths.fr/alea/static/e3c/2020/tex/e3c_2020_00_sujet49_3_cor.tex
\begin{document}
\noindent\textit{APMEP / CoopMaths. CC BY-SA. Présentation adaptée par Maths Explorer.}\par\medskip
\subsection*{1.}

On a donc \( u_2 = 200 + 5 = 205 \) et \( u_3 = u_2 + 5 = 205 + 5 = 210 \).

De même, on a :
\[
v_2 = 200 \times \left(1 + \dfrac{2}{100}\right) = 200 \times 1{,}02 = 204 \quad \text{et} \quad v_3 = v_2 \times 1{,}02 = 208{,}80 \text{ (en euros)}.
\]

\subsection*{2.}

\paragraph{a.}
\begin{center}
\begin{lstlisting}[style=python]
u = 200
v = 200
n = int(input("Saisir une valeur de n :"))
for i in range(1,n) :
    u = u + 5 
    v = v * 1.02
print("Pour n =",n,"on a","u =",u,"et v =",v)
\end{lstlisting}
\end{center}

\paragraph{b.} On a :
\[
u_4 = u_3 + 5 = 215 \quad \text{et} \quad v_4 = v_3 \times 1{,}02 \approx 110{,}98.
\]

\subsection*{3.}

Quel que soit le naturel \( n \), \( n \geq 1 \), on a :
\begin{itemize}
    \item \( u_{n+1} = u_n + 5 \) : la suite \( (u_n) \) est une suite arithmétique de premier terme 200 et de raison 5, donc \( u_n = 200 + 5(n - 1) \).
    \item \( v_{n+1} = v_n \times 1{,}02 \) : la suite \( (v_n) \) est une suite géométrique de premier terme 200 et de raison 1{,}02, donc \( v_n = 200 \times 1{,}02^{n-1} \).
\end{itemize}

\subsection*{4.}

3 ans correspondent à \( 3 \times 12 = 36 \) mois de loyers.

\begin{itemize}
    \item Avec le premier contrat :
    \[
    S_{36} = u_1 + u_2 + \dots + u_{36},
    \]
    que l'on peut écrire :
    \[
    S_{36} = 200 + 205 + \dots + 200 + 5 \times 35,
    \]
    ou encore :
    \[
    S_{36} = 200 + 5 \times 35 + \dots + 205 + 200,
    \]
    et en sommant membres à membres :
    \[
    2S_{36} = 36 \times (200 + 200 + 5 \times 35) = 36 \times 575 = 20700,
    \]
   	d'où \(S_{36} = 10350\) (euros).

    \item Avec le second contrat :
    \[
    T_{36} = v_1 + v_2 + \dots + v_{36},
    \]
    que l'on peut écrire :
    \[
    T_{36} = 200 + 200 \times 1{,}02 + \dots + 200 \times 1{,}02^{35} \quad (1),
    \]
    d'où, par produit par 1{,}02 :
    \[
    1{,}02 T_{36} = 200 \times 1{,}02 + 200 \times 1{,}02^2 + \dots + 200 \times 1{,}02^{36} \quad (2),
    \]
    et par différence \((2) - (1)\) :
    \[
    0{,}02 T_{36} = 200 \times 1{,}02^{36} - 200,
    \]
    d'où :
    \[
    T_{36} = \dfrac{200 \times 1{,}02^{36} - 200}{0{,}02} \approx 10398{,}87 \text{ (euros)}.
    \]
\end{itemize}

C'est donc le premier contrat qui reviendra le moins cher (de peu!).

\end{document}
